\documentclass[12pt, letterpaper]{article}

%-------Packages---------
\usepackage{amssymb,amsfonts}
\usepackage{mathptmx}
\usepackage{amsthm}
\usepackage{graphicx}
\usepackage[all,arc]{xy}
\usepackage[colorlinks=true, allcolors=blue]{hyperref}
\usepackage{enumerate}
\usepackage{bbm}
\usepackage{dsfont}
\usepackage{mathrsfs}
\usepackage{tikz-cd}
\usepackage{setspace}
\usepackage{import}
\usepackage[utf8]{inputenc}
\usepackage{hyperref}
\usepackage{tikz}
\usepackage{float}
\usepackage{mdframed}
\usepackage{fancyhdr}
\usepackage[nointegrals]{wasysym}

\usepackage[left=1in,top=1in,right=1in,bottom=1in]{geometry}

\usepackage{mathtools}
\usepackage{setspace}
\usepackage{cleveref}
\usepackage{multicol}

%--------Theorem Environments--------
%theoremstyle{plain} --- default
\newmdtheoremenv{theorem}{Theorem}[subsection]
\newmdtheoremenv{corollary}[theorem]{Corollary}
\newtheorem{proposition}[theorem]{Proposition}
\newmdtheoremenv{lemma}[theorem]{Lemma}
\newtheorem{conj}[theorem]{Conjecture}
\newtheorem{quest}[theorem]{Question}

\theoremstyle{definition}
\newmdtheoremenv{definition}[theorem]{Definition}
\newmdtheoremenv{definitions}[theorem]{Definitions}
\newtheorem{example}[theorem]{Example}
\newtheorem{algorithm}[theorem]{Algorithm}
\newtheorem{rmk}[theorem]{Remark}
\newtheorem{exercise}[theorem]{Exercise}

%----------- my commands -------------
\newcommand{\C}{\mathbb{C}}
\newcommand{\R}{\mathbb{R}}
\newcommand{\Q}{\mathbb{Q}}
\newcommand{\Z}{\mathbb{Z}}
\newcommand{\N}{\mathbb{N}}
\renewcommand{\P}{\mathbb{P}}
\newcommand{\contr}{\Rightarrow \Leftarrow}
\newcommand{\HRule}[1]{\rule{\linewidth}{#1}}
\newcommand{\s}{\(\,\)}
\renewcommand{\t}[1]{\text{#1}}
\newcommand{\ang}[1]{\left\langle #1 \right\rangle}

% Meta data
\setstretch{1.2}

\begin{document}

\pagestyle{fancy}
\fancyhf{}
\fancyhead[LOE]{\slshape{\textbf{Vincent Ferrara}}}
\fancyhead[ROE]{\thepage}

\subsection*{Problem 1}
$((c)\implies (a))$\\
Assume $(c)$\\
Let $\varphi: G \rightarrow H \times K$ s.t. $\varphi(g)=(\phi'(g),\psi(g))$\\
Homomorphism:\\
Since $\phi'$ and $\psi$ are homomorphisms $\varphi$ is a homomorphism.\\
Injective:\\
Let $g \in \t{ker}(\varphi)$\\
Thus $\varphi(g)=(\phi'(g),\psi(g))=(e,e)$ thus $g \in \t{ker}(\phi'),\,\t{ker}(\psi)$ 
thus since this is an exact sequence $\t{ker}(\psi)=\t{im}(\phi)$ thus since $\phi'\circ \phi=id_H \implies g=e$ thus ker$(\varphi)=\{e\}$\\
Surjective:
Let $(h,k) \in H \times K$, thus since $\sigma$ is surjective there exists a $g \in G$ s.t. $\psi(g)=k$ and for the $h$ consider $\phi(h)$\\
Thus $\varphi(\phi(h),g)=(\phi'(\phi(h)),\psi(g))=(h,k)$ thus surjective. Thus $G \cong H \times K$ thus proven in class since this is a direct product this sequence is split.\\
$((a)\implies (b))$\\
Assume $(a)$\\
Thus shown in class this means there exists a $K' \cong K$ s.t. $G=\phi(H)K$ and $\phi(H) \cap K'=1$\\
Consider $\psi\mid_{K'}: K' \rightarrow K$ is a homomorphism.\\
Injective:\\
Let $k' \in \t{ker}(\psi\mid_{K'})$ thus $\psi\mid_{K'}(k')=e \in \phi(H) \implies k'=e$ since $\phi(H) \cap K' = \{e\}$\\
Surjective:\\
For any $k \in K$, since $\psi$ is surjective, there exits $g \in G$ s.t. $\psi(g)=k$\\
Thus you can write $g=\phi(h)k'$ for $h \in H$ and $k' \in K'$\\
$k=\psi(g)=\psi(\phi(h)k')=\psi(\phi(h))\psi(k')=\psi(k')$\\
Thus $\psi\mid_{K'}$ is surjective.\\
Since $\psi\mid_{K'}$ is an isomorhpism take its inverse and define $\phi': K \rightarrow G$ s.t. $\psi'=(\psi\mid_{K'})^{-1}$ we can do thus since $K' \subset G$\\
$\psi'$ is a homomorphism since $(\psi\mid_{K'})^{-1}$ is a homomorphism.\\
Consider $\psi(\psi'(k))$ for $k \in K$ well as shown above $\psi'$ is actually only going to map to $K'$ and because of this we can think of $\psi=\psi\mid_{K'}$ thus since these are inverse of eachother\\
$\psi(\psi'(k))=k$ thus done\\
$((b)\implies (a))$\\
Assume $(b)$\\
Define $K' = \psi'(K)$\\
Let $x \in \phi(H) \cap K'$\\
Thus $x=\phi(h)=\psi'(k)$ for $h \in H$ and $k \in K$\\
$\psi(\phi(h))=\psi(\psi'(k)) \implies e=k \implies x=\psi'(e)=e$ thus $\phi(H) \cap K'=\{e\}$\\
Let $g \in G$\\
Define $k'=\psi'(\psi(g)) \in K'$\\
Consider $\psi(gk'^{-1})=\psi(g)\psi(k'^{-1})=\psi(g)\psi(k')^{-1}=\psi(g)\psi(\psi'(\psi(g)))^{-1}=\psi(g)\psi(g)^{-1}=e$\\
Thus $gk'^{-1} \in \t{ker}(\psi)=\t{im}(\phi)$ thus there is an $h \in H$ s.t. $\phi(h)=gk'^{-1}$ thus $g=\phi(h)k'$\\
Thus $G=\phi(H)K'$\\
Thus $G$ is split

\newpage
\subsection*{Problem 2}
For a group of order $8$ to be split in this sequence $G \cong C_2 \rtimes C_2 \times C_2$ thus there exists a $\phi: C_2 \times C_2 \rightarrow \t{Aut}(C_2)\cong\{e\}$
thus the only homomorphism $\phi(a)=id_{C_2\times C_2}$ s.t. $a \in C_2$ thus to be split $G \cong C_2 \times C_2 \times C_2$
\begin{enumerate}
    \item $C_8$\\
    If $N \cong C_2$, then $C_8 \diagup N \cong C_4$ thus $C_8$ cannot fit in this sequence.
    \item $C_4 \times C_2$\\
    $C_2 \cong \ang{(a^2,e)}$\\
    $C_4\times C_2 \diagup \ang{(a^2,e)} = \{\ang{(a^2,e)}, (a,e)\ang{(a^2,e)}, (e,b)\ang{(a^2,e)}, (a,e)\ang{(a^2,e)}\} \cong C_2 \times C_2$\\
    $e \rightarrow \ang{(a^2,e)} \rightarrow G \rightarrow G\diagup\ang{(a^2,e)} \rightarrow e$\\
    Thus this sequence fits and is not split from above.
    \item $C_2\times C_2 \times C_2$\\
    $C_2 \cong \ang{(a,e,e)}$\\
    $C_2\times C_2 \times C_2 \diagup \ang{(a,e,e)} = \{\ang{(a,e,e)}, (e,b,e)\ang{(a,e,e)}, (e,e,c)\ang{(a,e,e)}, (e,b,c)\ang{(a,e,e)}\} \cong C_2 \times C_2$\\
    $e \rightarrow \ang{(a,e,e)} \rightarrow G \rightarrow G\diagup\ang{(a,e,e)} \rightarrow e$\\
    Thus this sequence fits and is split by above.
    \item $D_4$
    $C_2 \cong \ang{x^2}$\\
    $D_4 \diagup \ang{x^2} = \{\ang{x^2}, x\ang{x^2}, y\ang{x^2}, xy\ang{x^2}\} \cong C_2 \times C_2$\\
    $e \rightarrow \ang{x^2} \rightarrow G \rightarrow G\diagup\ang{x^2} \rightarrow e$\\
    Thus this sequence fits and is not split from above.
    \item $Q$
    $C_2 \cong \ang{-1}$\\
    $Q \diagup \ang{-1} = \{\ang{-1}, i\ang{-1}, j\ang{-1}, k\ang{-1}\} \cong C_2 \times C_2$\\
    Thus this sequence fits and is not split from above.
\end{enumerate}

\newpage
\subsection*{Problem 3}
\hspace{0.48cm} If \( G \) is split, there must exist a subgroup \( D_2' \cong D_2 \) in \( G \). This implies we can center all reflections and rotations at a single point, as there is a subgroup isomorphic to \( D_2 \).

Let \( v \) be the smallest vector on the \( x \)-axis, and let \( w \) be the vector with the smallest height. Let \( z \) be the smallest vector in the lattice.

\begin{itemize}
    \item If \( z \) lies on the \( x \)-axis, then \( z = v \).
    \item If \( z \) lies on the \( y \)-axis, then \( z = w \).
\end{itemize}

Assuming without loss of generality that \( z \) is in the first quadrant, we can act on it by \( r \), so \( r(z) \) is in our lattice. Then \( r(z) + z \) lies on the \( x \)-axis. Since \( z \) is the smallest vector in the lattice, we have:
\[
x(r(z) + z) < 2x(v) \Rightarrow 2x(z) < 2x(v)
\]
If \( 2x(w) \neq x(v), 0 \), then \( z + r(z) - v \) would be a smaller vector than \( v \) on the \( x \)-axis, which is a contradiction.

Thus, if \( x(z) = 0 \), then \( z = w \). If \( x(z) = \frac{1}{2}x(v) \), then by similar reasoning, we find \( x(z) = x(w) \). Since \( w \) has the smallest height and \( z \) is the smallest vector, it must be \( w \). Therefore, our lattice is spanned by \( v \) and \( w \).

We now have two cases: whether \( w \) is vertical or not.

1. \textbf{If \( w \) is vertical}, this corresponds to case 9:
   \[
   G \cong \mathbb{Z} \times \mathbb{Z} \rtimes D_2 = \langle x, y, r, p_\pi \mid xy = yx, \, rp_\pi = p_\pi r, \, r^2 = e, \, p_\pi^2 = e, \, r x r = x, \, r y r = y^{-1}, \, p_\pi x p_\pi = x^{-1}, \, p_\pi y p_\pi = y^{-1} \rangle
   \]
   with
   \[
   \phi(r) = \begin{pmatrix} 1 & 0 \\ 0 & -1 \end{pmatrix}, \quad \phi(p_\pi) = \begin{pmatrix} -1 & 0 \\ 0 & -1 \end{pmatrix}
   \]

2. \textbf{If \( w \) is not vertical}, this corresponds to case 12:
   \[
   G \cong \mathbb{Z} \times \mathbb{Z} \rtimes D_2 = \langle x, y, r, p_\pi \mid xy = yx, \, rp_\pi = p_\pi r, \, r^2 = e, \, p_\pi^2 = e, \, r x r = x, \, r y r = x y^{-1}, \, p_\pi x p_\pi = x^{-1}, \, p_\pi y p_\pi = y^{-1} \rangle
   \]
   with
   \[
   \phi(r) = \begin{pmatrix} 1 & 1 \\ 0 & -1 \end{pmatrix}, \quad \phi(p_\pi) = \begin{pmatrix} -1 & 0 \\ 0 & -1 \end{pmatrix}
   \]
\phantom{text}\\
If \( G \) is not split, then we have two cases:

\begin{enumerate}
    \item If one reflection lifts to only a glide.
    \item If both reflections lift to only a glide.
\end{enumerate}

For the first case, if one reflection lifts to a glide, we make the glide parallel to the \( x \)-axis and the reflection parallel to the \( y \)-axis, centering the rotations at the origin.
Let the min glide (meaning moves the minimum distance) parallel to the \( x \)-axis be defined as \( t_{(a,h)}rt_{(0,-h)} = t_{(a,2h)}r \). Let a reflection parallel to the \( y \)-axis be \( t_{(b,0)}r_y t_{(-b,0)} = t_{(2b,0)}r_y \).

Now, consider:
\[
t_{(a,2h)}r t_{(a,2h)}r = t_{(a,2h)} t_{(a,-2h)} rr = t_{(2a,0)}
\]
Thus, we have a vector on the \( x \)-axis. We call the minimal vector on the \( x \)-axis \( v \). Assume \( v = (c,0) \) with \( c < a \). Then consider:
\[
t_{(-c,0)} t_{(a,2h)} r = t_{(a-c,2h)} r
\]
Since \( a - c < a \), this contradicts the minimal glide assumption, so \( v = (2a,0) \).\\
Similarly, we can find a vector on the \( y \)-axis by acting with \( r_y \), and then finding the minimal vector, which we call \( w \).\\
Let \( z \) be the smallest vector. If \( z \) lies on any axis, it must be \( v \) or \( w \). If the \( x \)- or \( y \)-coordinate of \( z \) is zero, then it is either \( v \) or \( w \), leaving us with \( z = (a, 0.5y(w)) \) using a similar argument as above. 

Consider the element:
\[
t_{-z} t_{(a,2h)} r = t_{(0,2h - 0.5y(w))} r
\]
Then:
\[
t_{(0,2h - 0.5y(w))} r t_{(0,2h - 0.5y(w))} r = t_{(0,2h - 0.5y(w))} t_{(0,-2h + 0.5y(w))} rr = \text{id}
\]
This is a reflection parallel to the horizontal axis, which is not allowed. Thus, this minimal vector \( z \) cannot exist. Therefore, the minimal vector is either \( v \) or \( w \), so these two span the lattice.\\
Thus:
\[
t_{(a,2h)} r = t_{\frac{1}{2}v + cw} r \in G \quad \text{for } c \in \mathbb{R}
\]
and
\[
t_{(2b,0)} r_y = t_{kv} r_y \in G \quad \text{for } k \in \mathbb{R}
\]
Now, consider:
\[
t_{kv} r_y t_{\frac{1}{2}v + cw} r = t_{(k - \frac{1}{2})v + cw} p_\pi \in G
\]
Thus:
\[
t_{(k - \frac{1}{2})v + cw} p_\pi p_\pi = t_{(k - \frac{1}{2})v + cw} \in G
\]
implying \( k - \frac{1}{2} \in \mathbb{Z} \) and \( c \in \mathbb{Z} \) since this translation is in my group, and we only have integer linear combinations of $v$ and $w$.\\
So:
\[
t_{-cw} t_{\frac{1}{2}v + cw} r = t_{\frac{1}{2}v} r = t_{(a,0)} r \in G
\]
Also, since \( k - \frac{1}{2} \in \mathbb{Z} \) and \( 2b = 2ak \Rightarrow b = ak \Rightarrow k = \frac{d}{2} \) for \( d \) odd, we have \( b = \frac{ad}{2} \). Setting \( d = 1 \), we find \( b = \frac{a}{2} \), so:
\[
t_{(a,0)} r_y \in G
\]
Thus, we have case 11:
\phantom{text}\\
$G = \langle \mu, y, r, p_\pi \mid \mu^2 y = y \mu^2, \, r^2 = p_\pi^2 = e, \, r \mu r = \mu^{-1}, \, r p_\pi = p_\pi r, \, r y r = y, \, \mu y \mu^{-1}, \, r \mu^{-1} = \mu r, \, p_\pi \mu^{-1} = \mu p_\pi, \, p_\pi \mu p_\pi = \mu^{-1}, \, p_\pi y p_\pi = y^{-1} \rangle$
\phantom{text}\\

If two reflections lift to glides, we pick the minimal glides \( t_{(a, 2c)} r \) and \( t_{(2d, b)} r_y \) parallel to their respective axes. The minimal vectors on the \( x \)-axis and \( y \)-axis, as in the previous proof, are \( v = (2a, 0) \) and \( w = (0, 2b) \), respectively.\\
Then:
\[
t_{(a,2c)} r = t_{\frac{1}{2}v + k w} r \quad \text{for } k \in \mathbb{R}
\]
\[
t_{(2d,b)} r_y = t_{j v + \frac{1}{2}w} r_y \quad \text{for } j \in \mathbb{R}
\]
Now consider:
\[
t_{j v + \frac{1}{2}w} r_y \, t_{\frac{1}{2}v + k w} r = t_{\left(\frac{1}{2} - j\right)v + \left(\frac{1}{2} + k\right)w} p_\pi \in G
\]
Thus:
\[
t_{\left(\frac{1}{2} - j\right)v + \left(\frac{1}{2} + k\right)w} p_\pi p_\pi = t_{\left(\frac{1}{2} - j\right)v + \left(\frac{1}{2} + k\right)w} \in G
\]
so we require \( \frac{1}{2} - j \in \mathbb{Z} \) and \( \frac{1}{2} + k \in \mathbb{Z} \).

Since \( k + \frac{1}{2} \in \mathbb{Z} \) and \( 2c = 2ak \), it follows that \( k = \frac{f}{2} \) for some odd integer \( f \). Choosing \( f = 1 \), we get \( c = \frac{a}{2} \). Similarly, \( d = \frac{b}{2} \), so \( t_{(a,0)} r \in G \) and \( t_{(0,b)} r_y \in G \).

Thus, we have case 10:
\phantom{text}\\
$G = \langle \mu, \gamma, p_\pi \mid \mu^2 \gamma^2 = \gamma^2 \mu^2, \, \gamma \mu^{-1} = \mu \gamma^{-1}, \, \mu p_\pi \mu^{-1} = \mu^2 p_\pi, \, \mu \gamma^{-1} = \gamma \mu^{-1}, \, \gamma p_\pi \gamma^{-1} = \gamma^2 p_\pi \rangle$
\phantom{text}\\

\newpage
\subsection*{Problem 4}
\hspace{0.48cm} If \( G \) lifts one reflection to a reflection s.t. a rotation of order 4 is on that line; thus, make the reflection like the $x$-axis, let this reflection be \( r \) with the center of the rotation \( p_{\pi/2} \) at the origin.

Let \( x \) be a vector. We can act on \( x \) by \( r \) to obtain something on the \( x \)-axis. Let \( v \) be the minimal vector on the \( x \)-axis.

Applying a rotation, we obtain a vector on the \( y \)-axis. Let \( w \) be the smallest vector on the \( y \)-axis.

Now, consider \( p_{\pi/2}(v) - w \); this vector would be shorter than \( w \), which is a contradiction. Thus, \( p_{\pi/2}(v) = w \), meaning they must be the same vector.

Let \( z \) be the smallest vector and, without loss of generality, assume it lies in the first quadrant. By a similar argument as above, we either conclude that \( z \) is on one of the axes or that \( z = \left(0.5 \cdot x(v), 0.5 \cdot x(v)\right) \).

If this is our minimal vector, we can rotate by \( 45^\circ \) to get a new coordinate system where the lattice is spanned by \( z \) and \( p_{\pi/2}(z) \) or by \( v \) and \( w \).

Thus, for any coord. system we will have a line of reflection with a center of an order 4 rotation.

Then, we can use $p_\pi/2$ on the reflection to get all reflections

Thus, this is a clear subgroup in $G$ that is isomorphic to $D_4$ thus we are in the split case

Thus, with a square lattice we have case 13:
\phantom{text}\\
$G \cong \mathbb{Z} \times \mathbb{Z} \rtimes D_4 = \langle x, y, r, p_{\pi/2} \mid xy = yx, \, r p_{\pi/2} = p_{\pi/2}^{-1} r, \, r^2 = e, \, p_{\pi/2}^4 = e, \, r x r = x, \, r y r = y^{-1}, \, p_{\pi/2} x p_{\pi/2}^{-1} = y, \, p_{\pi/2} y p_{\pi/2}^{-1} = x^{-1} \rangle
$\phantom{text}\\

For:
\[
\phi(r) = \begin{pmatrix} 1 & 0 \\ 0 & -1 \end{pmatrix}, \quad \phi(p_{\pi/2}) = \begin{pmatrix} 0 & -1 \\ 1 & 0 \end{pmatrix}
\]

If \( G \) does not lift any reflection to a reflection s.t. a rotation of order 4 is on that line.

Thus, this is the non split case because we cannot get 4 reflections intersection at one point which is needed for a subgroup isomorphic to $D_4$.\\

Thus, a reflection lifted to only glide, thus assume for contradiction that any glide is also not on a center of rotation of order 4.

Thus, 4 glides are on a 




\newpage
\subsection*{Problem 5}

\end{document}